% theory_optimal_capital.tex — Optimal Capital Structure for BTC Treasury Vehicles
% Extends Leland (1994) / Leland-Toft (1996) to multi-channel BTC acquisition

\section{Optimal Capital Structure}\label{sec:optimal_capital}

This section derives the optimal financing policy for a BTC treasury vehicle
that can issue common equity, convertible debt, and multiple series of
preferred stock.  We extend the classical Leland~(1994) and
Leland--Toft~(1996) endogenous capital structure framework to a setting where
the firm's sole asset is BTC and the objective is to maximize BTC per share
rather than firm value per se.

%% ====================================================================
\subsection{Channel-Specific Cost of BTC Acquisition}\label{subsec:channel_costs}

Each financing channel $c \in \mathcal{C} = \{\text{ATM},\,\text{conv},\,
\text{pfd-nc},\,\text{pfd-c},\,\text{ops}\}$ carries a distinct effective
cost of BTC acquisition, measured as the all-in cost per BTC acquired
relative to the spot price $S_t$.

\subsubsection{ATM Equity Issuance}

\begin{proposition}[Cost of ATM Issuance]\label{prop:cost_atm}
The effective cost of BTC acquisition via ATM equity issuance, expressed as
a dilution-adjusted cost per BTC, is
\begin{equation}\label{eq:k_atm}
  k_t^{\mathrm{ATM}}
  = S_t \left(1 - \frac{e^{\pi_t}}{\mathrm{ELE}_t}\right)^{-1}
    \cdot \frac{1}{e^{\pi_t} / \mathrm{ELE}_t},
\end{equation}
which simplifies to
\begin{equation}\label{eq:k_atm_simple}
  k_t^{\mathrm{ATM}} = S_t \cdot \frac{\mathrm{ELE}_t}{e^{\pi_t}}.
\end{equation}
The cost is below spot (i.e.\ accretive) when
$\pi_t > \log(\mathrm{ELE}_t)$, at spot when
$\pi_t = \log(\mathrm{ELE}_t)$, and above spot (dilutive) when
$\pi_t < \log(\mathrm{ELE}_t)$.
\end{proposition}

\begin{proof}
Issuing $\Delta N$ shares at price $P_t$ raises $P_t \Delta N$ in proceeds
and purchases $P_t \Delta N / S_t$ BTC.  The dilution cost to existing
shareholders per BTC acquired equals the reduction in BTC per share times
$N_t^{\mathrm{sh}}$:
\[
  \Delta B = B_t^{\mathrm{post}} - B_t
  = \frac{H_t + P_t \Delta N / S_t}{N_t^{\mathrm{sh}} + \Delta N} - B_t.
\]
In the marginal limit $\Delta N \to 0$, the effective cost per BTC is
$S_t \cdot (B_t S_t / P_t) = S_t \cdot (H_t S_t / E_t)
= S_t \cdot A_t / (e^{\pi_t}(A_t - L_t))
= S_t \cdot \mathrm{ELE}_t / e^{\pi_t}$.
This is less than $S_t$ iff $e^{\pi_t} > \mathrm{ELE}_t$, i.e.\
$\pi_t > \log(\mathrm{ELE}_t)$, consistent with the accretion condition
of Proposition~\ref{prop:accretion}.
\end{proof}

\begin{remark}
When $\pi_t \gg \log(\mathrm{ELE}_t)$, the effective cost approaches zero:
equity issuance at extreme premiums acquires BTC nearly ``for free'' in
per-share terms.  This is the economic engine of the flywheel at high premiums.
\end{remark}

\subsubsection{Convertible Debt}

\begin{proposition}[Cost of Convertible Debt Issuance]\label{prop:cost_conv}
The effective cost of BTC acquisition via convertible bond issuance with
face value $F$, coupon rate $\delta_d$, maturity $T_d$, and conversion ratio
$\phi_d$ is
\begin{equation}\label{eq:k_conv}
  k_t^{\mathrm{conv}}
  = S_t \left(1
    + \frac{\delta_d}{r} \bigl(1 - e^{-r T_d}\bigr)
    + p_{\mathrm{conv}}(t) \cdot \frac{\phi_d F \cdot \mathbb{E}_t[P_\tau]}{F}
    \right),
\end{equation}
where $p_{\mathrm{conv}}(t) = \mathbb{P}_t(\text{conversion before } T_d)$
is the risk-neutral conversion probability, $\tau$ is the conversion time,
and $\mathbb{E}_t[P_\tau]$ is the expected share price at conversion.
\end{proposition}

\begin{proof}
The all-in cost comprises three components: (i) the spot price $S_t$ paid
for BTC; (ii) the present value of coupon payments
$\delta_d F \cdot (1 - e^{-rT_d})/r$ per unit face value; and (iii) the
expected dilution cost at conversion, equal to the probability of
conversion times the equity value transferred to bondholders.  Each dollar
of face value purchases $1/S_t$ BTC, so the cost per BTC scales by $S_t$.
\end{proof}

\begin{remark}
Before conversion, the cost is the fixed coupon.  After conversion, the
ongoing coupon obligation vanishes but dilution is realized.  The two
regimes can be summarized:
\begin{equation}\label{eq:conv_regimes}
  k^{\mathrm{conv}} =
  \begin{cases}
    S_t\bigl(1 + \delta_d / r\bigr) & \text{if never converted (lower bound)}, \\
    S_t\bigl(1 + \phi_d \mathbb{E}[P_\tau] / S_t\bigr)
      & \text{if converted immediately (upper bound)}.
  \end{cases}
\end{equation}
Strategy's convertible notes have historically carried coupons of 0--0.875\%,
making them among the cheapest channels when conversion probability is low.
\end{remark}

\subsubsection{Non-Convertible Preferred Stock}

\begin{proposition}[Cost of Non-Convertible Preferred]\label{prop:cost_pfd_nc}
For a perpetual non-convertible preferred series with dividend rate
$\delta_p$ (e.g.\ STRF at $\delta_{\mathrm{STRF}} = 0.10$), the effective
cost of BTC acquisition is
\begin{equation}\label{eq:k_pfd_nc}
  k_t^{\mathrm{pfd\text{-}nc}}
  = S_t \left(1 + \frac{\delta_p}{r}\right),
\end{equation}
where $r$ is the discount rate.  This reflects the perpetual dividend
obligation as a lump-sum cost.
\end{proposition}

\begin{proof}
Each dollar of preferred proceeds purchases $1/S_t$ BTC.  The present value
of the perpetual dividend stream is $\delta_p / r$ per dollar of par.  The
total cost per BTC acquired is thus $S_t(1 + \delta_p / r)$.
\end{proof}

\begin{remark}
At $\delta_p = 0.10$ and $r = 0.05$, the effective cost multiplier is
$1 + 0.10/0.05 = 3\times$ spot, making non-convertible preferred the most
expensive channel.  This channel is only optimal when equity and convertible
markets are closed (e.g.\ during deep NAV discounts).  For STRF specifically,
the 10\% perpetual dividend creates a permanent cash flow obligation that
increases the distress boundary per Proposition~\ref{prop:distress}.
\end{remark}

\subsubsection{Convertible Preferred Stock}

\begin{proposition}[Cost of Convertible Preferred]\label{prop:cost_pfd_c}
For a convertible preferred series with dividend rate $\delta_c$ and
conversion ratio $\phi_c$ (e.g.\ STRK at $\delta_{\mathrm{STRK}} = 0.08$,
$\phi_c = 0.1$), the effective cost is
\begin{equation}\label{eq:k_pfd_c}
  k_t^{\mathrm{pfd\text{-}c}}
  = S_t \left(1
    + \frac{\delta_c}{r}\bigl(1 - p_{\mathrm{conv}}^{\mathrm{pfd}}(t)\bigr)
    + p_{\mathrm{conv}}^{\mathrm{pfd}}(t) \cdot
      \frac{\phi_c \mathbb{E}_t[P_\tau]}{S_t}
    \right),
\end{equation}
where $p_{\mathrm{conv}}^{\mathrm{pfd}}(t)$ is the probability of eventual
conversion.  If converted, the dividend obligation ceases but dilution is
realized; if never converted, the full perpetual dividend is borne.
\end{proposition}

\begin{proof}
The cost decomposes into: (i) spot BTC acquisition cost $S_t$; (ii)
expected present value of dividends paid until conversion
$(\delta_c/r)(1 - p_{\mathrm{conv}}^{\mathrm{pfd}})$, approximating the
dividend duration; and (iii) expected dilution at conversion.  The two
costs are mutually exclusive conditional on conversion, hence weighted
by conversion probability.
\end{proof}

%% ====================================================================
\subsection{Optimal Financing Rule}\label{subsec:optimal_rule}

We now derive the optimal channel selection as a function of the state
vector $(\pi_t, S_t, \mathrm{ELE}_t, \mathrm{DCR}_t)$.

\begin{definition}[Premium Thresholds]\label{def:thresholds}
Define the following critical premium thresholds:
\begin{align}
  \pi^*_{\mathrm{ATM}} &= \log(\mathrm{ELE}_t),
    \label{eq:pi_atm_thresh} \\
  \pi^*_{\mathrm{conv}} &= \log\!\left(\frac{k_t^{\mathrm{conv}}}{S_t}
    \cdot \frac{A_t - L_t}{A_t}\right),
    \label{eq:pi_conv_thresh} \\
  \pi^*_{\mathrm{pfd}} &= -\infty \quad (\text{always available, no premium
    requirement}).
    \label{eq:pi_pfd_thresh}
\end{align}
The ATM threshold $\pi^*_{\mathrm{ATM}}$ is the accretion boundary from
Proposition~\ref{prop:accretion}.  The convertible threshold
$\pi^*_{\mathrm{conv}}$ is the premium at which the convertible cost equals
the ATM cost.
\end{definition}

\begin{theorem}[Optimal Channel Selection]\label{thm:optimal_channel}
Suppose the firm minimizes the effective cost of BTC acquisition per share.
The optimal financing channel at time $t$ is determined by the premium
$\pi_t$ relative to the threshold values:
\begin{equation}\label{eq:optimal_channel}
  c^*(t) =
  \begin{cases}
    \mathrm{ATM} & \text{if } \pi_t > \pi^*_{\mathrm{ATM}}
      \text{ and ATM capacity available}, \\[4pt]
    \mathrm{conv} & \text{if } \pi^*_{\mathrm{conv}} < \pi_t
      \le \pi^*_{\mathrm{ATM}}
      \text{ and vol.\ sufficiently high}, \\[4pt]
    \mathrm{pfd\text{-}c} & \text{if } \pi_t \le \pi^*_{\mathrm{conv}}
      \text{ and DCR}_t > 1, \\[4pt]
    \mathrm{pfd\text{-}nc} & \text{if } \pi_t \le \pi^*_{\mathrm{conv}}
      \text{ and DCR}_t \le 1 \text{ (last resort)}.
  \end{cases}
\end{equation}
\end{theorem}

\begin{proof}
We establish that the cost ordering
$k^{\mathrm{ATM}} \le k^{\mathrm{conv}} \le k^{\mathrm{pfd\text{-}c}}
\le k^{\mathrm{pfd\text{-}nc}}$ holds when $\pi_t > \pi^*_{\mathrm{ATM}}$,
so ATM dominates.

\emph{Step 1.}  When $\pi_t > \log(\mathrm{ELE}_t)$, the ATM cost
$k^{\mathrm{ATM}} = S_t \cdot \mathrm{ELE}_t / e^{\pi_t} < S_t$, while all
other channels have cost $\ge S_t$.  Thus ATM strictly dominates.

\emph{Step 2.}  When $\pi_t \le \log(\mathrm{ELE}_t)$, ATM issuance is
dilutive ($k^{\mathrm{ATM}} \ge S_t$).  Convertible debt with a low coupon
$\delta_d$ may have $k^{\mathrm{conv}} < k^{\mathrm{ATM}}$.  The indifference
point is $\pi^*_{\mathrm{conv}}$, below which convertibles dominate ATM.

\emph{Step 3.}  When $\pi_t$ is too low for attractive convertible pricing
(conversion premium demands become extreme) or the equity is trading at or
below NAV, preferred issuance remains available since it does not depend on
the equity premium.  Convertible preferred (lower coupon) dominates
non-convertible preferred (higher coupon) when the firm can credibly
offer conversion rights, i.e.\ when $\mathrm{DCR}_t > 1$.

\emph{Step 4.}  Each channel is further subject to capacity constraints
(Section~\ref{subsec:constraints}), creating a lexicographic ordering where
the firm uses the cheapest available channel first, then moves down the
cost ladder as capacity is exhausted.
\end{proof}

\begin{corollary}[Strategy's Observed Behavior]\label{cor:observed}
Strategy's historical financing activity is broadly consistent with
Theorem~\ref{thm:optimal_channel}:
\begin{enumerate}
  \item During 2024--2025, with $\pi_t \in [0.5, 2.5]$ and
    $\log(\mathrm{ELE}_t) \approx 0.15$, the firm issued
    \$14.6B+ via ATM equity, the cheapest channel.
  \item Convertible debt (\$7.3B in 2024--2025) was issued when BTC
    volatility was elevated, enabling low-coupon structures
    ($\delta_d = 0\text{--}0.875\%$).
  \item Preferred stock (STRK, STRF, STRC, STRD, STRE) was issued in
    early 2026 as the premium compressed, consistent with the model's
    prediction that preferred becomes optimal at lower premiums.
\end{enumerate}
\end{corollary}

%% ====================================================================
\subsection{Weighted Average Cost of BTC Acquisition}\label{subsec:wacba_opt}

Building on Definition~\ref{def:wacba}, we now characterize WACBA as an
optimization target.

\begin{definition}[Instantaneous WACBA]\label{def:wacba_inst}
The instantaneous Weighted Average Cost of BTC Acquisition at time $t$ is
\begin{equation}\label{eq:wacba_inst}
  \mathrm{WACBA}_t
  = \sum_{c \in \mathcal{C}} w_c(t) \cdot k_c(\pi_t, S_t, \mathrm{ELE}_t),
\end{equation}
where $w_c(t) = dH_t^{(c)} / dH_t$ is the fraction of BTC acquired through
channel $c$ and $\sum_c w_c(t) = 1$.
\end{definition}

\begin{proposition}[WACBA Minimization]\label{prop:wacba_min}
Management's optimization problem is
\begin{equation}\label{eq:wacba_opt}
  \min_{\{w_c(t)\}_{c \in \mathcal{C}}} \mathrm{WACBA}_t
\end{equation}
subject to:
\begin{align}
  & w_c(t) \ge 0 \quad \forall\, c, \qquad \sum_c w_c(t) = 1,
    & &\text{(non-negativity, partition)} \label{eq:wacba_c1} \\
  & \int_0^T dH_t^{\mathrm{ATM}} \le \bar{V}_{\mathrm{ATM}} \cdot T,
    & &\text{(ATM market capacity)} \label{eq:wacba_c2} \\
  & \frac{D_s + D_c + P_{nc} + P_c}{A_t} \le \bar{\ell},
    & &\text{(leverage ceiling)} \label{eq:wacba_c3} \\
  & \Phi_t \le \bar{\Phi}(A_t, \mathrm{CF}_t^{\mathrm{ops}}),
    & &\text{(dividend coverage floor)} \label{eq:wacba_c4}
\end{align}
where $\bar{V}_{\mathrm{ATM}}$ is the daily ATM volume capacity,
$\bar{\ell}$ is the maximum leverage ratio, and $\bar{\Phi}$ is the
maximum sustainable dividend obligation.
\end{proposition}

\begin{proof}
The objective is linear in $\{w_c\}$ with convex constraints.  The solution
is a corner solution (bang-bang): use the cheapest available channel up to
its capacity, then the next cheapest.  This follows from the simplex
structure of the feasible set and linearity of the objective.  The capacity
constraints~\eqref{eq:wacba_c2}--\eqref{eq:wacba_c4} bind sequentially as
cheaper channels are exhausted.
\end{proof}

\begin{remark}
In practice, Strategy does not literally solve this program but follows a
heuristic that approximates it: aggressive ATM issuance at high premiums
(cheapest channel, largest capacity), opportunistic convertible issuance
during volatility spikes, and preferred issuance as a backstop when
equity channels are unattractive.
\end{remark}

%% ====================================================================
\subsection{Endogenous Leverage Dynamics}\label{subsec:endo_leverage}

A critical feature of the BTC treasury model is that leverage evolves
endogenously and procyclically, creating amplification dynamics akin to
those in \citet{brunnermeier2009market}.

\begin{theorem}[Procyclical Leverage]\label{thm:procyclical}
Define the leverage ratio $\ell_t = L_t / A_t$.  Under the extended model
dynamics, leverage is countercyclical in BTC price and procyclical in
financial fragility:
\begin{enumerate}
  \item \textbf{Bull market ($dS_t > 0$):}
    $A_t$ rises while $L_t$ is approximately constant, so $\ell_t$ falls.
    Simultaneously, $\pi_t$ tends to rise (positive BTC--premium correlation
    regime), enabling ATM issuance that purchases BTC without increasing
    liabilities.  This further reduces $\ell_t$.
  \item \textbf{Bear market ($dS_t < 0$):}
    $A_t$ falls while $L_t$ remains fixed (debt and preferred par values
    do not adjust), so $\ell_t$ rises.  The premium compresses
    ($\pi_t$ falls), shutting off the ATM channel.  If the firm issues
    preferred stock to acquire BTC, $L_t$ increases, further raising
    $\ell_t$.
\end{enumerate}
\end{theorem}

\begin{proof}
Differentiate $\ell_t = L_t / (H_t S_t)$:
\begin{equation}\label{eq:dldt}
  d\ell_t
  = -\frac{L_t}{A_t^2}\, dA_t + \frac{dL_t}{A_t}
  = -\ell_t \cdot \frac{dA_t}{A_t} + \frac{dL_t}{A_t}.
\end{equation}
In a bull market, $dA_t > 0$ and $dL_t \approx 0$ (ATM issuance does not
add liabilities), so $d\ell_t < 0$.  In a bear market, $dA_t < 0$ and
$dL_t \ge 0$ (preferred issuance adds to $L_t$), so $d\ell_t > 0$.

The feedback loop arises because rising leverage increases
$\mathrm{ELE}_t = 1/(1 - \ell_t)$, amplifying the equity response to
further BTC declines, which raises leverage further.
\end{proof}

\begin{proposition}[Leverage Spiral Bound]\label{prop:leverage_spiral}
The leverage ratio has an absorbing barrier at $\ell_t = 1$
(insolvency: $A_t = L_t$).  The speed of approach to this barrier
accelerates as $\ell_t \to 1$ because the elasticity
$\mathrm{ELE}_t = 1/(1-\ell_t) \to \infty$.  Specifically, the
expected time to distress from leverage level $\ell$ satisfies
\begin{equation}\label{eq:distress_time}
  \mathbb{E}[\tau_{\mathrm{distress}} \mid \ell_t = \ell]
  \;\propto\; (1 - \ell)^2 / \sigma_S^2
\end{equation}
for $\ell$ close to 1, under GBM dynamics for $S_t$ with constant
liabilities.
\end{proposition}

\begin{proof}
With $L_t = L$ constant and $A_t = H S_t$, the process
$X_t = \log(A_t / L) = \log(H S_t / L)$ follows a Brownian motion with
drift $\mu_S - \sigma_S^2/2$ and volatility $\sigma_S$.  Distress occurs
when $X_t = 0$.  Starting from $X_0 = -\log\ell$, the expected first-passage
time to zero scales as $(\log(1/\ell))^2 / \sigma_S^2
\approx (1-\ell)^2 / \sigma_S^2$ for $\ell$ near 1.
\end{proof}

\begin{remark}[Connection to Brunnermeier--Pedersen]
The procyclical leverage mechanism mirrors the margin spiral of
\citet{brunnermeier2009market}: declining asset values tighten funding
constraints, forcing sales (or preventing purchases) that further depress
values.  In the Strategy context, the ``margin'' is the premium $\pi_t$,
which determines access to the cheapest financing channel (ATM equity).
Premium compression in bear markets plays the role of margin calls,
forcing the firm to use more expensive channels (preferred stock) that
increase leverage, creating a self-reinforcing downward spiral.
\end{remark}

%% ====================================================================
\subsection{Leland Framework Extension}\label{subsec:leland}

We now formalize the optimal capital structure in the tradition of
\citet{leland1994} and \citet{leland1996}, adapted to the
BTC treasury setting.

\begin{assumption}[Leland--BTC Environment]\label{ass:leland_btc}
\begin{enumerate}
  \item BTC price follows $dS_t/S_t = \mu_S\,dt + \sigma_S\,dW_t$ under
    $\mathbb{P}$.
  \item The firm holds $H$ BTC (fixed for this analysis) with asset value
    $A_t = H S_t$.
  \item Liabilities consist of perpetual debt with coupon $c_D$ and
    perpetual preferred with dividend $c_P$.
  \item The firm defaults when $A_t$ hits an endogenous boundary
    $A^* = \underline{A}$, chosen optimally by equity holders.
  \item At default, a fraction $\alpha$ of asset value is lost to
    bankruptcy costs.
\end{enumerate}
\end{assumption}

\begin{theorem}[Optimal Default Boundary]\label{thm:default_boundary}
Under Assumption~\ref{ass:leland_btc}, the equity-value-maximizing default
boundary is
\begin{equation}\label{eq:A_star}
  A^* = \frac{\gamma}{\gamma + 1} \cdot \frac{c_D + c_P}{r},
\end{equation}
where $\gamma = (r - \mu_S + \sigma_S^2/2) / \sigma_S^2
+ \sqrt{((r - \mu_S + \sigma_S^2/2)/\sigma_S^2)^2 + 2r/\sigma_S^2}$
is the positive root of the characteristic equation
$\tfrac{1}{2}\sigma_S^2 \gamma(\gamma-1) + \mu_S \gamma - r = 0$.
\end{theorem}

\begin{proof}
Equity value satisfies $E(A) = A - (c_D + c_P)/r + C \cdot A^{-\gamma}$
for $A > A^*$, where $C$ is determined by the boundary conditions
$E(A^*) = 0$ (limited liability) and $E'(A^*) = 0$ (smooth pasting).
Smooth pasting yields
$1 - \gamma C (A^*)^{-\gamma-1} = 0$, and combined with
$A^* - (c_D+c_P)/r + C(A^*)^{-\gamma} = 0$, solving for $A^*$ gives
the stated result.
\end{proof}

\begin{corollary}[Optimal Coupon / Dividend Split]\label{cor:optimal_split}
The total payout $c = c_D + c_P$ that maximizes firm value
$V(A) = E(A) + D(A) + P(A)$ is
\begin{equation}\label{eq:optimal_c}
  c^* = r \cdot A_0 \cdot \frac{\gamma + 1}{\gamma}
    \cdot \left(\frac{\tau_{\mathrm{eff}}}{1 + \tau_{\mathrm{eff}}
    (1-\alpha)}\right)^{1/\gamma},
\end{equation}
where $\tau_{\mathrm{eff}}$ is the effective tax advantage of
debt over preferred (in a BTC treasury context, this may reflect
regulatory or accounting benefits rather than traditional tax shields).
The split between $c_D$ and $c_P$ is then determined by market access
constraints per Theorem~\ref{thm:optimal_channel}.
\end{corollary}

\begin{remark}
A key difference from classical Leland: the BTC treasury firm has no
operating cash flows to service coupons, so all payments must come from
(i) new issuance, (ii) BTC sales, or (iii) legacy software revenue.
This makes the default boundary more fragile than for operating companies,
as the firm cannot ``earn its way out'' of distress.
\end{remark}
